\documentclass[11pt,onecolumn]{IEEEtran}
\usepackage[a4paper]{geometry}
\usepackage{amsmath,amsfonts,amssymb}
\usepackage{color}
\usepackage{hyperref}

\topmargin      -12.0mm
\oddsidemargin  -7.0mm
\evensidemargin -7.0mm
\textheight     233.0mm
\textwidth      173.0mm
\columnsep        4.1mm
\parindent        1.5em
\headsep          6.3mm
\headheight        12pt
\lineskip           1pt
\normallineskip     1pt
\def\baselinestretch{1.1}

\begin{document}

\begin{center}
    {\bf\Large Response to the Reviewer Comments on Manuscript}\\[2mm]
    {\bf\Large ``Federated Wind Speed Forecasting in Finland: Geographic vs.\ Correlation Graphs with Synchronous FedGD''}\\[7mm]
    {\em\large corresponding author: Anonymous}\\[6.0mm]
    \today
\end{center}

\vspace{7mm}

We thank the reviewers and the teaching assistant for their constructive and detailed feedback.
We have addressed all raised issues in the revised manuscript.
Modifications in the revised manuscript are highlighted in red.

\vspace{2mm}

Major modifications include:

\vspace{2mm}

\begin{itemize}

    \item The report title was revised to convey the main idea (comparing graph structures),
          the application domain (wind speed forecasting in Finland), and the methods (FL with synchronous FedGD).

    \item All in-plot figure numbers were corrected to match their caption labels, center-aligned,
          and redundant overlapping titles were removed.

    \item The feature count was corrected from 7 to 6, the prediction label was formally defined
          with mathematical notation, and additional justification for the squared $\ell_2$ norm was added.

    \item The Future Work section was expanded from a short paragraph to a detailed bullet list
          with three concrete research directions.

\end{itemize}

In the following, we respond to each comment point by point.
Section, figure, and equation numbers in the quoted comments refer to the original manuscript;
numbers in our responses refer to the revised manuscript.

\vspace{5mm}


%%%%%%%%%%%%%%%%%%%%%%%%%%%%%%%%%%%%%%%%%%%%
\begin{center}
    {\bf\large Comments of Reviewer \#1}
    \vspace{1mm}
\end{center}
%%%%%%%%%%%%%%%%%%%%%%%%%%%%%%%%%%%%%%%%%%%%

\begin{enumerate}

    \item[1.1] {\em The report should include the list of FMI station IDs or names used (e.g., as a table or
          supplementary reference), since the $\geq$80\% completeness filter applied to an undated API call
          cannot reliably reproduce the exact 50 stations.}

          \vspace*{2mm}
          \begin{itemize}
              \item The data section (Section 4) was expanded to include the exact WFS stored query name
                    (\texttt{fmi::observations::weather::hourly::simple}), the bounding box coordinates,
                    all four variable names, the completeness threshold (80\% of 8\,784 hourly observations),
                    and the gap-filling rule, so the filtering procedure can be reproduced exactly.
          \end{itemize}
          \vspace*{4mm}

    \item[1.2] {\em There is an inconsistency where figure numbers embedded inside plots (e.g., ``Figure 4'',
          ``Figure 7'', ``Figure 8'') do not match the caption labels below them (``Fig.\ 2'', ``Fig.\ 5'',
          ``Fig.\ 6'') --- these should be cleaned up before the final submission.}

          \vspace*{2mm}
          \begin{itemize}
              \item All in-plot titles were corrected: Fig.\ 2 no longer reads ``Figure 4'',
                    Fig.\ 5 no longer reads ``Figure 7'', and Fig.\ 6 no longer reads ``Figure 8''.
                    All in-plot titles are now center-aligned.
          \end{itemize}
          \vspace*{4mm}

    \item[1.3] {\em Some of the texts are overlapping and affect the readability of plots (Fig.\ 6 / Figure 8).
          Since there are already captions below each figure, the captions inside the plots seem unnecessary
          and reduce legibility.}

          \vspace*{2mm}
          \begin{itemize}
              \item The redundant in-plot titles were removed from Fig.\ 6.
                    The figure caption below already conveys the same information.
          \end{itemize}
          \vspace*{4mm}

\end{enumerate}


%%%%%%%%%%%%%%%%%%%%%%%%%%%%%%%%%%%%%%%%%%%%
\begin{center}
    {\bf\large Comments of Reviewer \#2}
    \vspace{1mm}
\end{center}
%%%%%%%%%%%%%%%%%%%%%%%%%%%%%%%%%%%%%%%%%%%%

\begin{enumerate}

    \item[2.1] {\em The report discusses the use of 7 features, but it only explains 6:
          3 lagged wind speeds (t-1, t-2, t-3), 1 lagged atmospheric pressure, 1 lagged temperature,
          1 lagged humidity.}

          \vspace*{2mm}
          \begin{itemize}
              \item The feature count was corrected to 6 throughout the manuscript.
                    The weight vector $\mathbf{w}^{(i)} \in \mathbb{R}^7$ has 7 components because the
                    feature vector is augmented with a bias term, not because there are 7 features.
                    This distinction is now made explicit in Section 2.
          \end{itemize}
          \vspace*{4mm}

    \item[2.2] {\em Feature normalization is repeated unnecessarily in the Problem Formulation section.}

          \vspace*{2mm}
          \begin{itemize}
              \item The redundant normalization description was removed from the Problem Formulation section.
                    Feature and label normalization are now described exactly once, in Section 2,
                    immediately after the feature vector definition.
          \end{itemize}
          \vspace*{4mm}

    \item[2.3] {\em GTVMin (and FedGD) is discussed in Section 2 before being explained in Section 3.
          It is also not cited. The reader should not be assumed to be previously familiar with GTVMin.}

          \vspace*{2mm}
          \begin{itemize}
              \item A forward reference ``(defined in Section~3)'' and a citation to \cite{SarcheshmehPour2023}
                    were added at the first mention of GTVMin in Section~2.
                    FedGD is similarly cited at first use.
          \end{itemize}
          \vspace*{4mm}

    \item[2.4] {\em For the Geographic k-NN system, the edge weight kernel explanation is a bit unclear:
          it uses a Gaussian kernel (not mentioned or cited); the report makes it sound like all of
          Equation~2 is the Haversine distance; the Haversine formula is not shown or cited.}

          \vspace*{2mm}
          \begin{itemize}
              \item The edge weight formula (Eq.~2) is now explicitly described as a Gaussian kernel.
              \item The variable $d_{ij}$ in Eq.~2 is now clearly identified as the Haversine distance
                    between stations $i$ and $j$, so it is clear that the Haversine formula gives $d_{ij}$
                    and not the full equation.
              \item A citation to the Haversine formula \cite{Sinnott1984} was added.
          \end{itemize}
          \vspace*{4mm}

    \item[2.5] {\em Some subsections would improve the report's readability.}

          \vspace*{2mm}
          \begin{itemize}
              \item Subsections were not added. The spconf IEEE template is designed for short papers
                    and does not support subsections well. The report is already at the page limit,
                    so adding subsection headers would require removing content.
          \end{itemize}
          \vspace*{4mm}

    \item[2.6] {\em The figures have inconsistent number labeling in their titles and in the report.}

          \vspace*{2mm}
          \begin{itemize}
              \item All in-plot figure numbers were corrected to match their caption labels,
                    as described in our response to comment~1.2.
          \end{itemize}
          \vspace*{4mm}

\end{enumerate}


%%%%%%%%%%%%%%%%%%%%%%%%%%%%%%%%%%%%%%%%%%%%
\begin{center}
    {\bf\large Comments of Teaching Assistant (Alexander Pavlyuk)}
    \vspace{1mm}
\end{center}
%%%%%%%%%%%%%%%%%%%%%%%%%%%%%%%%%%%%%%%%%%%%

\begin{enumerate}

    \item[TA.1] {\em Criterion 1: The title is not informative. Please indicate (i) the main idea (comparing
          graph structures), (ii) the application domain (wind prediction), and (iii) the methods
          (FL, Synchronous FedGD).}

          \vspace*{2mm}
          \begin{itemize}
              \item The title was changed to ``Federated Wind Speed Forecasting in Finland:
                    Geographic vs.\ Correlation Graphs with Synchronous FedGD''.
                    It now conveys (i) the comparison of graph structures (Geographic vs.\ Correlation Graphs),
                    (ii) the application domain (wind speed forecasting in Finland), and
                    (iii) the method (Synchronous FedGD).
          \end{itemize}
          \vspace*{4mm}

    \item[TA.2] {\em Criterion 6: The label was not explicitly defined. Mathematical notation should be added
          to improve readability.}

          \vspace*{2mm}
          \begin{itemize}
              \item The prediction label is now formally defined as $y^{(i)}_t = ws^{(i)}_{t+1}$
                    (next-hour wind speed at station $i$) at the start of Section~2.
              \item The full feature vector is now written out as a math expression with all six
                    components named and their physical units stated.
          \end{itemize}
          \vspace*{4mm}

    \item[TA.3] {\em Criterion 9: More arguments in favor of the squared $\ell_2$ norm can be made.}

          \vspace*{2mm}
          \begin{itemize}
              \item Three explicit arguments for the squared $\ell_2$ norm were added in Section~3:
                    (1) it is smooth and differentiable everywhere, giving a closed-form gradient
                    that pairs directly with the FedGD update rule;
                    (2) it penalizes large parameter disagreements quadratically, encouraging
                    tight agreement between neighbors rather than merely bounded spread;
                    (3) adding a strictly convex penalty to the sum of convex local losses makes
                    the full GTVMin objective strongly convex for $\alpha > 0$, guaranteeing a
                    unique global minimizer.
          \end{itemize}
          \vspace*{4mm}

    \item[TA.4] {\em Criterion 14: Figure title/numbering issues; y-axis label out of bounds on Figure 3;
          in-plot title overlaps; station coordinates used instead of names.}

          \vspace*{2mm}
          \begin{itemize}
              \item All in-plot figure numbers were corrected to match their caption labels,
                    as described in our response to comment~1.2.
              \item The y-axis label that was clipped out of bounds on Fig.~3 was repositioned.
              \item Station coordinates used to label the best and worst performing stations
                    were replaced with their FMI station names for readability.
          \end{itemize}
          \vspace*{4mm}

    \item[TA.5] {\em Criterion 16: I strongly suggest including more detailed potential improvements
          and presenting them in a list.}

          \vspace*{2mm}
          \begin{itemize}
              \item The Future Work section was rewritten from a short paragraph into a bullet list
                    with three concrete directions: (1) more expressive local models (MLP, LSTM),
                    (2) sparser data regimes where graph regularization provides more benefit, and
                    (3) a wind-flow-aware directed graph based on prevailing wind directions.
          \end{itemize}
          \vspace*{4mm}

\end{enumerate}

\vspace{5mm}

\end{document}
